\begin{center}
\large\textcolor{Red}{\textbf{Introduction to the Breviary}}
\end{center}

\begin{multicols}{2}

\ubsection{Introduction to the Breviary}{introduction-breviarium}

% \vspace{+1em}

% {\color{Red} \ubsubsection{Lost traditions}{}}

{\color{Red} \ubsubsection{I. Content \& Structural Overview}{content-and-structural-overview}}

As a \emph{Totum} Breviary, this volume contains everything needed for a traveling Friar Preacher to recite daily, as well as the Little Office of the Blessed Virgin, along with rubrics on the communal obligation of the Office of the Dead, among other rubrics. I have included the Necrology of the Masters-General, as they are written in the Prototype, namely until 1345, and included the rubric on how to execute the termination of collects.

This volume aims to be a diplomatic transcription of the manuscript, adding accent marks strictly within the \emph{Psalterium}, as well as asterisks for Psalmody, Invitatories, and Responsories. Doxologies in Hymns have been expanded where abbreviated. Tone markings in the lessons have been omitted. Medieval orthography has been kept, while adopting u's and v's for modern conveniences.

% As I have kept faithful to the layout and order of the source manuscript, the rubrics are scattered in a variety of places, from the rubrics excerpted from the Collectarium (including the Kalendar), to the rubrics concerning the Temporal Cycle before Advent, to Rubrics concerning Sundays Per Annum placed after Epiphany, to rubrics concerning the Sanctoral placed before the Sanctoral Propers, placed in April are the Commons for Paschaltide, and rubrics at the very end governing the supplemental offices.

As I have kept faithful to the layout and order of the source manuscript, the rubrics are scattered in a variety of places. The Breviary is laid out accordingly:

\begin{itemize}[nosep,leftmargin=10pt]
	\item Front Matter - The Preface and Introduction
	\item \emph{Martyrologium} - Contains the necrology of Masters-General and Anniversaries of the order, then the rubric governing the Gospel at the \emph{Preciosa}, or when the Constitution is read instead.
	\item \emph{Collectarium} - Contains the Kalendar of feasts, and rubrics elsewhere not found in the \emph{Breviarium}, such as the Blessing for Compline.
	\item \emph{Psalterium} - Contains the Psalter, the ferial Antiphons, and the Daily Office of the Blessed Virgin.
	\item \emph{Breviarium} - Contains the following:
	\begin{itemize}[nosep,leftmargin=0pt]
		\item \emph{Rubrice Temporum} - Rubrics regarding the Sunday Office
		\item \emph{Proprium de Tempore} - Rubrics and Propers for the Season.
		\item \emph{Dedicatione Ecclesie} - Rubrics and Propers for the Dedication of a Church.
		\item \emph{Rubrice Sanctorum} - Rubrics for the Feasts of Saints.
		\item \emph{Proprium Sanctorum} - Propers for the Feasts of Saints. The Common of Saints for Paschaltide is also found here.
		\item \emph{Commune Sanctorum} - Commons for the Saints, first the Office, then how to commemorate them, including Blessed Dominic, then the lessons for Matins. Then, rubrics for the Saturday Conventual Office of the Blessed Virgin, the Gradual Psalms, the \emph{Salve} and \emph{Ave} after Compline, the Office for the Dead, the \emph{Preciosa} Office, and lastly, the Blessings at Matins.
		\end{itemize}
	\item Indices - Indices for the Psalms and Canticles, and the Hymns.
\end{itemize}

% {\color{Red} \ubsubsection{IV. Comparative Guide}{comparative-guide}}
{\color{Red} \ubsubsection{II. Practical Rubrics \& \ Private Recitation}{practical-rubrics}}

Humbert himself wrote a commentary on the Constitutions of the Dominican Order, in which he commented on the \emph{Ordinarium} as well, which covers partly on the Divine Office, translated by Fr. Albert Judy, O.P., PhD, and available at: \emph{www.domcentral.org/study/}. Several rubrics are not written in the Breviary proper, and custom must be sought, which are explained in this Commentary. For example, we read at \S 17.32, rules regarding private recitation:
\begin{itemize}[nosep,leftmargin=8pt]
	\item The weekly nine-lesson Office of the Dead is not omitted,
	\item Antiphons are not intoned, but only said after the psalm.
	\item The Office is said loudly and clearly, not mumbled or ``between the teeth''. Care should also be taken to pronounce the words distinctly and unhurriedly, nor should the mind or body be distracted.
	\item The same applies to dress, that a friar may be dressed well before saying the Office.
	\item Prostrations may be made, as well as inclinations. In such, the beginning of the hours, the Chapter, Collects, versicles, responsories, and blessings are said standing.
	\item Psalms 148-150 and the New Testament Canticles are said standing.
	\item The \emph{Pater noster} is never said sitting, which includes the \emph{Credo} and \emph{Ave}.
\end{itemize}


	\ubsubsection{On the Matins Lessons}{matins-lessons}
The Matins lessons presented in the \emph{Breviarium} are not the same as the lessons as written in the \emph{Lectionarium}. In this Breviary, lessons are to be repeated from the Common, where the rubrics state one is free to divide them to avoid weariness, except for a handful of feasts where the lessons are proper, and are written so. The reason for this is twofold: to alleviate the burden of copying a Breviary, and secondly, to place the emphasis on the Psalms at Matins, with the lessons serving their purpose as breaks from Psalmody, rather than the unfortunate mindset of lessons being the main focal point at Matins.

	\ubsubsection{The \emph{Pater} and \emph{Credo} Before the Hours}{pater-and-credo}
Before the hours of Matins and Prime, there is to be recited a \emph{Pater} and \emph{Credo}, as the rubrics remind us on Maundy Thursday and on All Souls' Day. Before other hours, except for Compline, a \emph{Pater} only. Likewise, a \emph{Pater} is said, after the \emph{Fidelium}, if another hour does not follow.


	\ubsubsection{On the Compline Blessing}{compline-blessing}
As friars were never to be absent from Compline, the Blessing for Compline is given by the hebdomadarian and hence not copied in the \emph{Breviarium}. Later Dominican Breviaries include the blessing, and so, in private, the wording is changed from \emph{vos} to \emph{nos}.

	\ubsubsection{The \emph{Salve} after Compline}{salve-procession}
After the \emph{Benedictio} at Compline, the Compline of the Daily Office of the Virgin is recited, then the \emph{Salve} or \emph{Ave} Antiphon is sung with its versicle and Collect \emph{Concede nos}. After which, the \emph{Fidelium anime} is said, followed by the \emph{Pater noster} and \emph{Credo}.


	\ubsubsection{The \emph{Pater noster}}{pater}
The \emph{Pater noster} is written in full only in the \emph{Missale Conventuale} and the \emph{Missale Minorum Altarium}, and not in the preceding books of the Prototype. Here it is per the \emph{Missale Conventuale}.
% \begin{multicols}{2}
\fullpsalm{pater}
% \end{multicols}

	% \ubsubsection{On Humbert's Commentary on the Constitutions}{commentary}


	% \ubsubsection{On the Layout of the Breviary}{breviary-layout}


	\ubsubsection{III. Overview of Psalm Distribution}{psalm-distribution}

Note that the Psalms are printed sequentially and not as they are to be prescribed per the ferial Office. However, the order in which they appear nearly matches when they are to be said, with a few exceptions:

\begin{itemize}[nosep,leftmargin=8pt]
	\item At Sunday Matins, they begin from 1 until 20, excepting Psalms 4, \psalm{4} and 5, \psalm{5}, which are said at Compline and Monday Lauds, respectively.
	\item Psalms 21-25 are recited at Sunday Prime during Septuagesimatide until Easter exclusive.
	\item The Psalms of Monday Matins are entirely from 26 to 37.
	\item The Psalms of Tuesday Matins are from 38 to 51, skipping Psalm 50, \psalm{50}, which is reserved for ferial Lauds.
	\item The Psalms of Wednesday Matins are sequentially from Psalm 52 to 67, excepting Psalms 53, \psalm{53} which is said at Prime daily, and Psalms 62, \psalm{62} and 66, \psalm{66}, which are recited every day at Lauds as a single Psalm, and Psalm 64, \psalm{64}, which is said later at Lauds.
	\item The Psalms of Thursday Matins are said entirely from 68 to 79.
	\item The Psalms of Friday Matins are from 80 to 96, leaving out Psalm 89, \psalm{89} for Thursday Lauds, 90, \psalm{90} at Compline, 91, \psalm{91} for Saturday Lauds, 92, \psalm{92} for Sunday and festal Lauds, and 94, \psalm{94} is recited only on Epiphany day.
	\item The Psalms of Saturday Matins are said entirely from 97 to 108.
	\item Psalms 109 to 113 are then recited at Sunday Vespers,
	\item Psalms 114 to 120 are then recited at Monday Vespers, excepting Psalms 117, \psalm{117} and 118, \psalm{118}, the former of which is reserved for Sunday Lauds during Septuagesimatide until Easter exclusive, and the latter is recited daily at the minor hours.
	\item Psalms 121 to 125 are then recited at Tuesday Vespers,
	\item Psalms 126 to 130 are then recited at Wednesday Vespers,
	\item Psalms 131 to 136 are then recited at Thursday Vespers, excepting Psalm 133, \psalm{133} for Compline.
	\item Psalms 137 to 141 are then recited at Friday Vespers,
	\item Psalms 143 to 147 are then recited at Saturday Vespers, placing Psalm 142, \psalm{142} for Friday Lauds, and
	\item Psalms 148-150 are repeated at Lauds daily throughout the year, as a single Psalm.
\end{itemize}

{\color{Red} \ubsubsection{IV. Comparative Guide}{comparative-guide}}

The Dominican Breviary will be very much familiar to anyone who has recited the 1962 Roman Office on I Class feasts, with seven canonical hours, and the nocturnal hours of Matins. As on I Class feasts, at the minor hours are recited Psalms 53 and 118, and at Compline, which is more or less fixed, are said the Sunday Psalms of 4, 90, and 133, with the addition of the first six verses of Psalm 30, and a rich variety of Antiphons.

At Matins, Sundays have 18 Psalms, ferias have 12, and feasts outside of Paschaltide have 9 Psalms, but in Paschaltide, 3 Psalms are said daily. At Lauds, 7 Psalms are said plus the Old Testament canticle, with the 3rd and 5th slot reserved daily for Psalms 62 \& 66, and Psalms 148-150 daily.

At Prime, outside of Sundays between Septuagesima and Palm Sunday, it is fixed to Psalms 53 and the first four stanzas of 118. On those excepted Sundays, Psalms 21-25 and 117 are added. The \emph{Quicumque} is recited nearly every Sunday. Those accustomed to the Roman Office will notice the Chapter Office of Prime is decoupled from Prime, and is instead be recited in choir after Lauds or Prime, depending on the Prior. After the Martyrology is read, the \emph{Lectio Brevis} that follows is not that of the Season or the Nones Chapter, but either the beginning of the Gospel of the day, or the beginning of the Constitutions of the Order. The remaining minor hours are the same as the Roman Office, finishing the remainder of Psalm 118.

At Vespers, there are differing rubrics compared to the Roman Office on which Psalms and Antiphons were to be recited, depending on the rank of the feast. Antiphons are never doubled, except on Totum Duplex feasts, at the \emph{Magnificat} and \emph{Benedictus} Antiphons.


{\color{Red} \ubsubsection{V. Navigation \& \ Bookmarks}{how-to-navigate}}
Since this Breviary was printed via Print-on-Demand, there are no ribbons, and so users are encouraged to use plenty of prayer cards and sticky tabs to keep their place. The following are recommended:

\noindent Bookmarks/Prayer Cards:
\begin{itemize}[nosep,leftmargin=8pt]
	\item At the Psalter, Temporal, Sanctoral, Common, and Common Lessons
\end{itemize}


\noindent Sticky Tabs:
\begin{itemize}[nosep,leftmargin=8pt]
	\item Canticles/\emph{Te Deum}/\emph{Venite}
	\item Minor Hours (1st Sunday After the Octave of Epiphany)
	\item Blessings/\emph{Preciosa}
	\item Daily Commemoration of St. Dominic
	\item \emph{Confiteor}/\emph{Salve}
	\item Compline Blessing (\emph{Collectarium})
\end{itemize}


{\color{Red} \ubsubsection{VI. Project Resources and \emph{Ordo}}{community-and-open-source-collaboration}}

It is recommended for the novice to acquire the \emph{Ordo} and Kalendar, available as a companion to this Breviary, at \textit{https://breviariumpredicatorum.pages.dev/\#ordo} as a guide to the Rubrics and Kalendar for each year, with pointers to which Feast, Psalms, Antiphons, Responsories, and Commemorations are said on each day. It also includes an English version of the rubrics in this Breviary.

This Breviary is an open-source project. Despite rigorous proofreading against the source manuscript, typos or other grammatical errors may persist in this first edition. If you find errors, or if you wish to contribute to the underlying LaTeX template, please submit an issue or pull request at: \emph{https://www.Github.com/AlanL399/BreviaryTypesetting/}.

Additionally, emails can be sent to \emph{BreviaryTypesetting@Outlook.com} if there are any questions, typos, or comments.





\end{multicols}

\begin{multicols}{2}

% {\color{Red} \ubsubsection{Test}{test-foreward}}

% \selectcolormodel{gray}

% {\color{Red} \ubsubsection{Test 2}{test2-foreward}}



% \textcolor{black!100}{the five boxing wizards jump quickly 100}
% \par \textcolor{black!90}{the five boxing wizards jump quickly 90}
% \par \textcolor{black!80}{the five boxing wizards jump quickly 80}
% \par \textcolor{black!70}{the five boxing wizards jump quickly 70}
% \par \textcolor{black!60}{the five boxing wizards jump quickly 60}
% \par \textcolor{black!50}{the five boxing wizards jump quickly 50}


% \lettrine[lines=2]{\zallmancaps \color{Red} R}{ubric} that has sections of text \ul{with underline} and \emph{italic text.} And more text. \ul{This is also some unique text that will be tested upon with the eyes, so don't be afraid to read on.} \emph{Here is some more text that is italicized, so please read on to test the reader. \ul{And we underline it as well, for fun's sake, to add more text to read.}} Red
% \lettrine[lines=2]{\zallmancaps \color{Blue} R}{ubric} that has sections of text \ul{with underline} and \emph{italic text.} And more text. \ul{This is also some unique text that will be tested upon with the eyes, so don't be afraid to read on.} \emph{Here is some more text that is italicized, so please read on to test the reader. \ul{And we underline it as well, for fun's sake, to add more text to read.}} Blue


% \textcolor{black!100}{waltz bad nymph for quick jigs vex 100}
% \par \textcolor{black!90}{waltz bad nymph for quick jigs vex 90}
% \par \textcolor{black!80}{waltz bad nymph for quick jigs vex 80}
% \par \textcolor{black!70}{waltz bad nymph for quick jigs vex 70}
% \par \textcolor{black!60}{waltz bad nymph for quick jigs vex 60}
% \par \textcolor{black!50}{waltz bad nymph for quick jigs vex 50}

% \begin{responsory-breve}
% {Sample Responsory With Text * And A Second Part.}{Plus The Versicle.}
% \end{responsory-breve}%

% \newline Otherwise, there are no italicized text, \ul{only underlined text.} \emph{So, please, dear reader, make the right choice.}

% \grecross \grealtcross \gothRbar \gothVbar \ \greheightstar \ \gresixstar \ 
% \grecross \grealtcross \gothRbar \gothVbar \ \greheightstar \ \gresixstar

% \color{Green}

% \grecross \grealtcross \gothRbar \gothVbar \ \greheightstar \ \gresixstar \ 
% \grecross \grealtcross \gothRbar \gothVbar \ \greheightstar \ \gresixstar

% \color{Red}

% \grecross \grealtcross \gothRbar \gothVbar \ \greheightstar \ \gresixstar \ 
% \grecross \grealtcross \gothRbar \gothVbar \ \greheightstar \ \gresixstar

% \color{Blue}

% \grecross \grealtcross \gothRbar \gothVbar \ \greheightstar \ \gresixstar \ 
% \grecross \grealtcross \gothRbar \gothVbar \ \greheightstar \ \gresixstar

% \color{Dandelion}

% \grecross \grealtcross \gothRbar \gothVbar \ \greheightstar \ \gresixstar \ 
% \grecross \grealtcross \gothRbar \gothVbar \ \greheightstar \ \gresixstar

% Test more colors, dvipsnames

% \color{Plum}

% \grecross \grealtcross \gothRbar \gothVbar \ \greheightstar \ \gresixstar \ 
% \grecross \grealtcross \gothRbar \gothVbar \ \greheightstar \ \gresixstar

% \color{MidnightBlue}

% \grecross \grealtcross \gothRbar \gothVbar \ \greheightstar \ \gresixstar \ 
% \grecross \grealtcross \gothRbar \gothVbar \ \greheightstar \ \gresixstar

% \color{BrickRed}

% \grecross \grealtcross \gothRbar \gothVbar \ \greheightstar \ \gresixstar \ 
% \grecross \grealtcross \gothRbar \gothVbar \ \greheightstar \ \gresixstar

% \color{SpringGreen}

% \grecross \grealtcross \gothRbar \gothVbar \ \greheightstar \ \gresixstar \ 
% \grecross \grealtcross \gothRbar \gothVbar \ \greheightstar \ \gresixstar

% \color{Mahogany}

% \grecross \grealtcross \gothRbar \gothVbar \ \greheightstar \ \gresixstar \ 
% \grecross \grealtcross \gothRbar \gothVbar \ \greheightstar \ \gresixstar

% \color{Maroon}

% \grecross \grealtcross \gothRbar \gothVbar \ \greheightstar \ \gresixstar \ 
% \grecross \grealtcross \gothRbar \gothVbar \ \greheightstar \ \gresixstar

% \color{Tan}

% \grecross \grealtcross \gothRbar \gothVbar \ \greheightstar \ \gresixstar \ 
% \grecross \grealtcross \gothRbar \gothVbar \ \greheightstar \ \gresixstar

% \color{OrangeRed}

% \grecross \grealtcross \gothRbar \gothVbar \ \greheightstar \ \gresixstar \ 
% \grecross \grealtcross \gothRbar \gothVbar \ \greheightstar \ \gresixstar

% \color{ForestGreen}

% \grecross \grealtcross \gothRbar \gothVbar \ \greheightstar \ \gresixstar \ 
% \grecross \grealtcross \gothRbar \gothVbar \ \greheightstar \ \gresixstar

% \color{PineGreen}

% \grecross \grealtcross \gothRbar \gothVbar \ \greheightstar \ \gresixstar \ 
% \grecross \grealtcross \gothRbar \gothVbar \ \greheightstar \ \gresixstar

% \color{RoyalBlue}

% \grecross \grealtcross \gothRbar \gothVbar \ \greheightstar \ \gresixstar \ 
% \grecross \grealtcross \gothRbar \gothVbar \ \greheightstar \ \gresixstar

% \color{Thistle}

% \grecross \grealtcross \gothRbar \gothVbar \ \greheightstar \ \gresixstar \ 
% \grecross \grealtcross \gothRbar \gothVbar \ \greheightstar \ \gresixstar

% {\color{Red} \ubsubsection{Test}{-introduction}}
% {\color{Red} \ubsubsection{Test}{-introduction}}
% {\color{Red} \ubsubsection{Test}{-introduction}}


% \lettrine[lines=2]{\zallmancaps \color{Red} O}{}
% \lettrine[lines=2]{\zallmancaps \color{Blue} I}{n}
% \lettrine[lines=2]{\zallmancaps \color{Red} O}{}
% \lettrine[lines=2]{\zallmancaps \color{Blue} I}{n}


\end{multicols}


\begin{center}
\vspace{+2em}
% \fontspec{Fraktur-Schmuck}\fontsize{48}{52}\selectfont B
% \newline
% \fontspec{Fraktur-Schmuck}\fontsize{48}{52}\selectfont C
% \newline
\fontspec{Fraktur-Schmuck}\fontsize{48}{52}\selectfont D
% \newline
% \fontspec{Fraktur-Schmuck}\fontsize{48}{52}\selectfont H
% \newline
% \fontspec{Fraktur-Schmuck}\fontsize{48}{52}\selectfont I
\end{center}


	% \ubsubsection{}{}
	% \ubsubsection{}{}
	% \ubsubsection{}{}
	% \ubsubsection{}{}
	% \ubsubsection{}{}
	% \ubsubsection{}{}
	% \ubsubsection{}{}